\documentclass{article}
\usepackage{graphicx}
\usepackage{booktabs, tabularx, amsmath, amssymb}
\newcolumntype{Y}[1]{>{\raggedright\arraybackslash\hsize=#1\hsize}X}

\title{ArtifactSweep}
\author{ahmedlamidi}
\date{September 2026}

\begin{document}
\maketitle

\section{Introduction}

\begin{table}[htbp]
\centering
\footnotesize
\renewcommand{\arraystretch}{1.3}
\setlength{\tabcolsep}{4pt}
\caption{Synthetic degradation models and real-world evaluation data.}
\label{tab:degradations}
\begin{tabularx}{\linewidth}{@{}>{\raggedright\arraybackslash}p{1.9cm} Y{1.2} Y{0.8}@{}}
\toprule
\textbf{Degradation} & \textbf{Forward model} & \textbf{Real-world data} \\
\midrule
Low-dose &
$y_i \sim \mathrm{Poisson}\big(\alpha b_0\, e^{-[Ax]_i}\big) + \mathcal{N}(0,\sigma_e^2)$ &
Mayo LDCT-PD; AAPM--Mayo 2016; 2DeteCT \\
\addlinespace
Sparse-view &
$\tilde p = \mathcal{P}_{\Omega}(Ax)$ \newline
$\Omega = \{\theta_{k\Delta}\}_{k=0}^{N_v-1} \subset \Theta$ &
Subsampled raw: Mayo LDCT-PD, Walnut CBCT, 2DeteCT \\
\addlinespace
Limited-angle &
$\tilde p = \mathcal{P}_{\Omega}(Ax)$ \newline
$\Omega = [\theta_0,\ \theta_0 + \theta_{\max}]$ &
Cropped raw: Mayo LDCT-PD, Walnut CBCT, 2DeteCT \\
\addlinespace
Beam hardening &
$y_i = \int \eta(E)\, e^{-a_i^\top \mu(E)}\, dE$ \newline
$\tilde p = -\ln\big(y / \textstyle\int \eta\big) \approx \sum_k c_k\, p^k$ &
2DeteCT (beam-hardening mode) \\
\addlinespace
Metal / starvation &
$\mu(E) = (1-M)\odot\mu_t(E) + M\odot\mu_m(E)$ \newline
$\bar y_i = \int \eta(E)\, e^{-a_i^\top \mu(E)}\, dE$ \newline
$y_i \sim \mathrm{Poisson}(\bar y_i + s_i) + \mathcal{N}(0,\sigma_e^2)$ &
CLINIC-metal; SpineWeb (no GT); AAPM CT-MAR 2023 (simulated) \\
\addlinespace
Motion &
$\tilde p(\theta) = A_\theta\big(\mathcal{T}_{\phi(\theta)}\, x\big)$ &
No paired benchmark; CQ500 (no GT) \\
\addlinespace
Ring / detector &
$\tilde p_{\theta,d} = g_d\, p_{\theta,d} + o_d$ \newline
$g_d = 0 \ \ \forall d \in \mathcal{D}_{\text{dead}}$ &
TomoBank (real rings) \\
\bottomrule
\end{tabularx}

\vspace{6pt}
\begin{minipage}{\linewidth}
\footnotesize
\textit{Notation and parameters.}
$x$ is the clean attenuation image, $A$ the system matrix with rows $a_i$, and $p = Ax$ the clean line integrals;
degraded sinograms are obtained as $\tilde p = -\ln(y/b)$, where $b$ is the unattenuated flux of each ray.
In the low-dose model, $b_0$ is the blank-scan photon count per detector bin, $\alpha \in (0,1]$ the dose fraction
relative to full dose, and $\sigma_e$ the standard deviation of electronic readout noise.
For sparse-view and limited-angle acquisitions, $\Theta$ is the full angular set and $\mathcal{P}_\Omega$ keeps only
the views in $\Omega$: sparse-view keeps $N_v$ evenly spaced views with step $\Delta = |\Theta|/N_v$ (typically 18--90 views),
whereas limited-angle keeps a contiguous arc of length $\theta_{\max}$ starting at $\theta_0$, with
$\theta_{\max} < \pi$ for parallel beam and $\theta_{\max} < \pi + 2\gamma$ for a fan beam of half-angle $\gamma$.
Beam hardening replaces the monoenergetic model with an integral over the source spectrum $\eta(E)$, set by tube
voltage and filtration, and energy-dependent attenuation $\mu(E)$; the polynomial with coefficients $c_k$ is the
common single-material approximation.
For metal, $M$ is the binary metal mask, $\mu_t(E)$ and $\mu_m(E)$ the tissue and metal attenuation, $\bar y_i$ the expected polychromatic count, and $s_i$ additive
scatter; photon starvation arises when $b_0 e^{-a_i^\top\mu}$ approaches zero along rays through metal, and the
clipping applied before the logarithm must be reported.
For motion, $A_\theta$ is the projector at view $\theta$ and $\mathcal{T}_{\phi(\theta)}$ warps the object with
view-dependent parameters $\phi(\theta)$: a rigid rotation and translation $(R_\theta, t_\theta)$ or a deformable field.
For ring artifacts, $d$ indexes detector elements, $g_d \sim 1 + \mathcal{N}(0,\sigma_g^2)$ is the per-element gain
error, $o_d$ the offset error, and $\mathcal{D}_{\text{dead}}$ the set of dead elements.
\end{minipage}
\end{table}

\end{document}
